% -*- coding: utf-8 -*-
% ============================================================
%  Template de slides — UFS Beamer
%  Baseado em: https://github.com/ayoshiaki/UFS_Beamer
%
%  Compilar localmente:
%    make          (usa o Makefile na raiz do repositório)
%  Ou manualmente:
%    TEXINPUTS=".:$(git rev-parse --show-toplevel)/theme/:" pdflatex slides.tex
% ============================================================

\documentclass[aspectratio=169]{beamer}

% --- Tema UFS --------------------------------------------------
% Os arquivos .sty e imagens estão em /theme/ na raiz do repo.
% O TEXINPUTS é configurado pelo Makefile / GitHub Actions.
\usetheme{UFS}

% --- Pacotes ---------------------------------------------------
\usepackage[utf8]{inputenc}
\usepackage[brazil]{babel}
\usepackage{amsmath,amssymb}
\usepackage{graphicx}
\usepackage{booktabs}
\usepackage{xcolor}
\usepackage{tikz}
\usetikzlibrary{positioning, shapes.multipart, arrows.meta,
                shadows, fit, chains,
                decorations.pathreplacing, shapes.geometric, calc}

% --- Configuração de listagens de código -----------------------
\usepackage{listings}
\definecolor{codegreen} {rgb}{0,   0.6, 0  }
\definecolor{codegray}  {rgb}{0.5, 0.5, 0.5}
\definecolor{codepurple}{rgb}{0.58,0,   0.82}
\definecolor{backcolour}{rgb}{0.95,0.95,0.92}

\lstdefinestyle{ufs}{
  backgroundcolor  = \color{backcolour},
  commentstyle     = \color{codegreen},
  keywordstyle     = \color{magenta},
  numberstyle      = \tiny\color{codegray},
  stringstyle      = \color{codepurple},
  basicstyle       = \ttfamily\tiny,
  breaklines       = true,
  captionpos       = b,
  numbers          = left,
  numbersep        = 5pt,
  frame            = single,
  rulecolor        = \color{gray},
  extendedchars    = true,
  literate =
    {ê}{{\^e}}1 {à}{{\`a}}1 {ú}{{\'u}}1
    {õ}{{\~o}}1 {ó}{{\'o}}1 {í}{{\'i}}1
    {á}{{\'a}}1 {ã}{{\~a}}1 {é}{{\'e}}1
    {ç}{{\c{c}}}1
}
\lstset{style=ufs}

% ===============================================================
%  METADADOS DA AULA — edite aqui
% ===============================================================
\title{Título da Aula}
\subtitle{Nome da Disciplina}
\author[Prof. André <andreyoshiaki@dcomp.ufs.br>]{Prof. Dr. André Yoshiaki Kashiwabara}
\institute{Departamento de Computação\\Universidade Federal de Sergipe}
\date{\today}

% Logo do departamento (opcional — remova a linha se não usar)
% \departmentlogo{logo_dcomp.png}

% ===============================================================
\begin{document}

% --- Capa ------------------------------------------------------
\begin{frame}
  \titlepage
\end{frame}

% --- Sumário ---------------------------------------------------
\begin{frame}{Sumário}
  \tableofcontents
\end{frame}

% ===============================================================
\section{Introdução}
% ===============================================================

\begin{frame}{Motivação}
  \begin{itemize}
    \item Por que estudar este tópico?
    \item Aplicações práticas
    \item O que veremos nesta aula
  \end{itemize}
\end{frame}

\begin{frame}{Objetivos}
  Ao final desta aula, o aluno será capaz de:
  \begin{enumerate}
    \item Objetivo 1
    \item Objetivo 2
    \item Objetivo 3
  \end{enumerate}
\end{frame}

% ===============================================================
\section{Desenvolvimento}
% ===============================================================

\begin{frame}{Conceito Principal}
  \begin{block}{Definição}
    Coloque aqui a definição formal do conceito.
  \end{block}

  \medskip
  \begin{exampleblock}{Exemplo}
    Ilustração prática do conceito.
  \end{exampleblock}
\end{frame}

\begin{frame}{Exemplo de Código}
\begin{lstlisting}[language=Python, caption={Exemplo em Python}]
def exemplo(n):
    """Docstring em português."""
    for i in range(n):
        print(f"Iteração {i}")
\end{lstlisting}
\end{frame}

\begin{frame}{Diagrama com TikZ}
  \begin{center}
  \begin{tikzpicture}[
    node distance=2cm,
    box/.style={draw, rounded corners, minimum width=2cm,
                minimum height=0.8cm, fill=UFSBlue!15}
  ]
    \node[box] (A) {Entrada};
    \node[box, right of=A] (B) {Processo};
    \node[box, right of=B] (C) {Saída};
    \draw[->, thick] (A) -- (B);
    \draw[->, thick] (B) -- (C);
  \end{tikzpicture}
  \end{center}
\end{frame}

% ===============================================================
\section{Exercícios}
% ===============================================================

\begin{frame}{Exercício Proposto}
  \begin{alertblock}{Exercício}
    Enunciado do exercício proposto.
  \end{alertblock}

  \medskip
  \textbf{Dica:} Dica ou orientação para resolução.
\end{frame}

% ===============================================================
\section{Conclusão}
% ===============================================================

\begin{frame}{Resumo}
  \begin{columns}
    \column{0.5\textwidth}
      \textbf{O que vimos:}
      \begin{itemize}
        \item Tópico 1
        \item Tópico 2
      \end{itemize}
    \column{0.5\textwidth}
      \textbf{Próxima aula:}
      \begin{itemize}
        \item Tópico A
        \item Tópico B
      \end{itemize}
  \end{columns}
\end{frame}

\begin{frame}{Referências}
  \begin{thebibliography}{9}
    \bibitem{ref1}
      Autor, A.
      \newblock \textit{Título do Livro}.
      \newblock Editora, Ano.
  \end{thebibliography}
\end{frame}

\end{document}
